\subsectionaddtolist{الف}
موجودیت یا Entity نماینده‌ی یک داده‌ی واقعی در سیستم و پایگاه داده است. 
این کلاس‌ها معمولا مستقیما با جدول‌های دیتابیس در ارتباط هستند و دارای \lr{ID} می‌باشند. 
همچنین ممکن است شامل منطق و رفتار مرتبط با دامنه‌ی مسئله باشند.

مثال:
\begin{latin}
\begin{verbatim}
class User {
    Long id;
    String name;
    String password;
}
\end{verbatim}
\end{latin}

\lr{DTO} یا \lr{Data Transfer Object} شیئی است که فقط برای انتقال داده بین لایه‌های مختلف سیستم استفاده می‌شود. 
\lr{DTO} معمولا منطق خاصی ندارد و فقط شامل اطلاعات موردنیاز است.

مثال:
\begin{latin}
\begin{verbatim}
class CreateUserDto {
    String name;
    String password;
}
\end{verbatim}
\end{latin}

تفاوت اصلی این است که \lr{Entity} نماینده‌ی داده‌ی اصلی سیستم و مرتبط با دیتابیس است؛
اما \lr{DTO} فقط برای انتقال داده و کاهش وابستگی یا افزایش امنیت استفاده می‌شود.

\subsectionaddtolist{ب}

اصل \lr{CCP} بیان می‌کند که کلاس‌هایی که به دلیل مشابهی تغییر می‌کنند، بهتر است در یک پکیج یا ماژول قرار بگیرند.

به بیان ساده، اگر چند کلاس معمولا هم‌زمان تغییر می‌کنند، باید کنار یکدیگر باشند تا مدیریت تغییرات ساده‌تر شود.

برای مثال، کلاس‌های مربوط به فاکتور مانند:
   \lr{Invoice,
  InvoiceService,
 InvoiceValidator}
بهتر است در یک پکیج مشترک قرار بگیرند؛
زیرا تغییر قوانین فاکتور احتمالا باعث تغییر همه‌ی آن‌ها می‌شود.

مزیت این اصل:
\begin{itemize}
    \item کاهش اثر تغییرات روی بخش‌های دیگر
    \item افزایش نگهداری‌پذیری پروژه
    \item ساده‌تر شدن توسعه و تست
\end{itemize}

\subsectionaddtolist{ج}
Decoupling به معنی کاهش وابستگی بین بخش‌های مختلف سیستم است. 
هرچه اجزای سیستم مستقل‌تر باشند، توسعه و نگهداری ساده‌تر خواهد بود.
منظور از سطح Deployment این است که اجزای سیستم چگونه منتشر و اجرا می‌شوند.
در یک سیستم کاملا Decoupled ، هر بخش می‌تواند به صورت مستقل Deploy شود.
برای مثال در معماری Microservice
 سرویس کاربران و  سرویس پرداخت و سرویس سفارش
هرکدام جداگانه اجرا و منتشر می‌شوند.
اما در معماری Monolith ، معمولا کل سیستم باید با هم Deploy شود.

مزایای Deployment مستقل:
\begin{itemize}
    \item سرعت بیشتر در انتشار تغییرات
    \item کاهش وابستگی بین نسخه‌ها
    \item امکان توسعه‌ی مستقل توسط تیم‌ها
    \item کاهش ریسک هنگام انتشار تغییرات
\end{itemize}